%% kapa: exam=1 round=23 year=2012 period=1 form=A subject=civil_law subject_ko=민법 seq=40 number=40
%% source: Q-Net 감정평가사 기출문제 (cst00309) artlSeq=5052578 file=raw/1008/5052578/2012년 제23회 감정평가사 제1차시험 기출문제(1교시 A형).pdf
%% conversion: kordoc 4.24.51 -> tools/gen1.py (auto); review=proofread-sonnet; flags=none
\begin{question}{40}
부합에 관한 설명으로 옳지 않은 것은?
\begin{choices}
  \item 동산이 부동산에 부합하여 동산의 소유권이 소멸한 때에는 그 동산을 목적으로 한 다른 권리도 소멸한다.
  \item 부합으로 인하여 권리를 상실하는 등 손해를 입은 자는 부당이득 규정에 의하여 보상을 청구할 수 있다.
  \item 건물로서의 독립성이 없어서 기존 건물의 구성부분이 된 경우에는 타인이 권원에 의해 부속시킨 것이라도 기존 건물에 부합한다.
  \item 저당권설정 후 저당목적물에 부합된 물건에는 저당권의 효력이 미치지 않는다.
  \item 건물 증축에 있어서 부합 여부는 증축 부분의 객관적 상태 외에 소유자의 의사를 고려하여 결정한다.
\end{choices}
\end{question}
\answer{4}
